\chapter{Related Work}
\label{sec:related-work}

This chapter reviews prior work in three areas relevant to the proposed approach: model predictive control for autonomous racing (Section~\ref{sec:rw-mpc}), kinodynamics modeling (Section~\ref{sec:rw-kinodynamics}), and vision foundation models for robotics (Section~\ref{sec:rw-vision}).

\section{Model Predictive Control for Autonomous Racing}
\label{sec:rw-mpc}

Model Predictive Control (MPC) has become a dominant paradigm for autonomous racing and aggressive vehicle navigation due to its ability to handle complex nonlinear dynamics, enforce state and input constraints, and optimize multi-objective cost functions over a receding horizon \cite{borrelli2012predictive, liniger2015optimization}. Traditional nonlinear MPC formulations rely heavily on gradient-based optimization techniques, such as sequential quadratic programming or interior-point methods \cite{diehl2002real, amrit2011fast}. These methods have been widely deployed using simplified kinematic or dynamic bicycle models, offering fast local convergence for standard driving maneuvers \cite{kong2015kinematic}. However, gradient-based solvers require continuous, twice-differentiable dynamics models and are highly susceptible to local minima \cite{verschueren2021acados}. Furthermore, at the extreme limits of handling, where tire saturation causes highly nonlinear and non-convex transitions, these methods often struggle to maintain real-time feasibility and computational stability \cite{subosits2019racetrack}.

To handle model uncertainties and disturbances, robust control frameworks such as Tube-based MPC have been introduced \cite{mayne2005robust}. While Tube-based MPC provides strict constraint satisfaction by optimizing over a nominal trajectory and bounding deviations within an invariant tube \cite{villanueva2017robust}, it tends to yield overly conservative behaviors. This conservatism is counterproductive in racing scenarios where operating near the friction limits is strictly required \cite{goh2019toward}.

Consequently, sampling-based MPC methods have gained significant traction as a gradient-free alternative. Algorithms such as the Cross-Entropy Method and Model Predictive Path Integral (MPPI) control generate thousands of candidate input sequences, propagate them through forward simulation, and compute the optimal control law via cost-weighted averaging \cite{williams2017information, wagener2019online}. MPPI, rooted in stochastic optimal control and information-theoretic principles, treats the underlying dynamics model as a completely black-box simulator \cite{williams2016aggressive}. This property makes MPPI well-suited for integration with complex, discontinuous, or neural network-based dynamics models that cannot be easily linearized \cite{drews2019vision, pinneri2021sample}.

Beyond single-execution optimization, Iterative learning control paradigms have been integrated with MPC to leverage historical driving data \cite{bristow2006survey}. Learning Model Predictive Control (LMPC) constructs a safe set and a terminal value function from prior feasible trajectories, systematically driving the vehicle to reduce lap times in subsequent iterations while maintaining recursive feasibility \cite{rosolia2017learning, rosolia2019learning}. Recent literature has extended this to sampling-based domains; for instance, Safe Information-Theoretic LMPC (SIT-LMPC) combines the parallelizable, gradient-free nature of MPPI with the iterative performance improvement and safety guarantees of LMPC \cite{yin2022safe, pravitra2020information}. Nevertheless, the performance of both gradient-based and sampling-based iterative controllers remains limited by the fidelity and adaptability of the forward predictive dynamics model \cite{kabzan2019learning}.

\section{Kinodynamics Modeling}
\label{sec:rw-kinodynamics}

Accurate forward simulation of vehicle states, particularly during high-slip maneuvers, requires kinodynamics models that account for inertial effects, actuation delays, and the highly nonlinear interactions between the tires and varying road surfaces \cite{rajamani2011vehicle}. Historically, autonomous navigation relied heavily on first-principles physics models equipped with empirical tire formulations, such as the Pacejka Magic Formula or Brush tire models \cite{pacejka2012tire, svendenius2005brush}. While these analytical models offer strict physical interpretability, their core assumption relies on steady-state tire behavior and uniform surface conditions. Calibration requires exhaustive offline system identification, and the fixed parameterizations fail entirely when the vehicle encounters unmodeled suspension dynamics, aerodynamic changes, or spatially varying terrain \cite{brunner2017sim, kaufmann2020deep}.

To circumvent the limitations of analytical parameterization, data-driven alternative architectures have been widely explored \cite{hou2020data}. Early attempts utilized simple Multi-Layer Perceptrons (MLPs) to map current states and actions to state derivatives \cite{punjani2015deep}. To better capture the temporal dependencies and actuation delays inherent in mechanical systems, recurrent architectures such as Long Short-Term Memory (LSTM) networks and Temporal Convolutional Networks (TCNs) have been employed to process historical observation windows \cite{lenz2015deep, d2020recurrent}. More recently, continuous-time formulations like Neural Ordinary Differential Equations (Neural ODEs) have been proposed to model underlying physical processes more naturally, effectively decoupling the dynamics from the sampling frequency of the sensors \cite{chen2018neural, finzi2020simplifying}. Despite their expressivity, purely deep-learning-based dynamics models often suffer from compounding integration errors over long prediction horizons and demonstrate poor generalization outside of their training distribution \cite{asadi2019combating}.

To balance physical priors with the flexibility of machine learning, hybrid or gray-box modeling has emerged as a robust compromise \cite{ajgl2020gray}. In this paradigm, a nominal first-principles model captures the dominant rigid-body kinematics, while a data-driven module, such as a Gaussian Process (GP) or a residual neural network, is trained to predict the unmodeled dynamics including aerodynamic drag or transient tire slip \cite{kabzan2019learning, ostafew2016robust}. GPs provide the added benefit of bounded uncertainty estimation, which is critical for robust control, though they scale poorly with large datasets \cite{hewing2020learning}.

However, a fundamental limitation persists across conventional analytical, purely learned, and hybrid models: they are predominantly reactive, relying entirely on proprioceptive sensors such as IMUs and wheel speed encoders \cite{lee2020learning}. When traversing heterogeneous environments, proprioception-only models can only register a change in surface friction after the vehicle has already experienced slip. This inherent latency prevents the controller from executing proactive, anticipatory maneuvers before engaging a hazardous or low-friction surface.

\section{Vision Foundation Models for Robotics}
\label{sec:rw-vision}

Proactive control in unstructured environments necessitates rich exteroceptive perception capable of anticipating physical interactions before they occur \cite{corke2023robotics}. Historically, visually guided navigation systems relied on explicit feature engineering, such as using Support Vector Machines on texture descriptors, to estimate a discrete surface class or a scalar friction coefficient \cite{broggi2013terrain, stano2022terrain}. Subsequently, end-to-end behavioral cloning and task-specific convolutional neural networks sought to map raw image pixels directly to steering and throttle commands \cite{bojarski2016end, codevilla2018end}. While conceptually elegant, these end-to-end policies are notoriously difficult to debug, require massive amounts of domain-specific human demonstration data, and rarely generalize across different vehicle platforms or lighting conditions \cite{muller2018driving}.

Self-supervised Vision Foundation Models (VFMs) have transformed visual perception for robotics \cite{bommasani2021opportunities}. Models trained via contrastive learning, such as CLIP and MoCo, or masked image modeling like MAE on internet-scale datasets, have demonstrated remarkable zero-shot generalization capabilities \cite{radford2021learning, he2020momentum, he2022masked}. Among these, DINO and its successor DINOv2 have proven particularly effective for dense robotic tasks \cite{caron2021emerging, oquab2023dinov2}. Because DINOv2 is optimized via self-distillation with no explicit labels, its intermediate attention maps and patch-level embeddings naturally encode deep semantic properties of the scene, such as texture, geometry, and material composition, without collapsing into task-specific biases \cite{amir2021deep}.

In the domain of field robotics and off-road navigation, VFMs have been actively leveraged to construct visual risk networks and traversability maps \cite{frey2023locomotion}. For instance, recent works have utilized frozen foundation model features to distinguish between navigable terrain, deep mud, and obstacles without fine-tuning, vastly outperforming supervised baselines in zero-shot deployment scenarios \cite{kahn2021badgr, seo2023scenepic}. However, integrating these high-dimensional semantic representations directly into vehicle dynamics remains an open problem.

The standard approach to utilizing visual information in vehicle control has been a modular, decoupled pipeline: a vision system segments the image or classifies the terrain type, assigns a predefined friction parameter based on a lookup table, and passes this updated parameter to the MPC's physics model \cite{srinivasan2020end, nitsche2022fast}. This explicit pipeline is brittle; it discards the nuanced, high-dimensional contextual information contained in the original embedding, propagates classification errors directly to the controller, and fails when encountering out-of-vocabulary surfaces that do not fit into predefined categories \cite{maturko2022vision}.

Among available VFMs, CLIP \cite{radford2021learning} and DINOv2 \cite{oquab2023dinov2} represent two complementary paradigms. CLIP aligns visual and language representations through contrastive pre-training on image-text pairs, producing embeddings that excel at semantic category-level tasks such as zero-shot classification. DINOv2, trained via self-distillation without language supervision, yields patch-level embeddings that encode finer-grained visual properties such as texture, material, and geometry. For dynamics adaptation, where the relevant signal is surface appearance rather than semantic category, DINOv2 embeddings provide a richer and more continuous representation of the visual context.

A practical challenge remains: both CLIP and DINOv2 produce high-dimensional embeddings (typically 384--768 dimensions) that are expensive to propagate through the dynamics model at every step of a sampling-based controller. PIVOT addresses this by processing visual features through a parameter network once per control cycle, producing a compact parameter vector that the lightweight dynamics model reuses across all rollout evaluations. This architecture decouples visual perception cost from the sampling and horizon dimensions of the controller, bypasses explicit terrain classification, and allows the dynamics to adapt continuously to the current operating environment.